\documentclass{article}
\usepackage[utf8]{inputenc}
\usepackage{graphicx}
\usepackage{hyperref}
\usepackage{listings}
\usepackage{xcolor}

% GATEHUB: command stubs are auto-injected on compile — no manual \newcommand needed.

\title{Learning Universe Compile Sample}
\author{THE GATEHUB}
\date{\today}

\begin{document}

\maketitle

\learninguniverse{
title={Sample Course},
description={Verifies all LMS commands compile without undefined control sequence errors.},
difficulty={Beginner},
estimatedHours={2},
skills={Python,Quizzes},
category={Computer Science}
}

\track{
title={Track 1},
description={First track},

\module{
title={Module 1},
description={Intro module},

\lesson{
title={Lesson 1: Basics}
}

\overviewmarkdown={
This lesson covers fundamentals. Use plain text here — no markdown hash headers.
}

\note{text={Remember to save your work.}}

\theory{title={Core Idea}, body={Theory content goes here.}}

\tip{text={Try the practice block before the quiz.}}

\warning{text={Do not skip the checkpoint.}}

\summary{text={You learned the basics of the DSL.}}

\keypoints{text={Commands compile, Parser works, PDF generates}}

\image{path={https://via.placeholder.com/400x200.png}, caption={Sample diagram}}

\video{type={youtube}, url={https://www.youtube.com/watch?v=dQw4w9WgXcQ}, title={Intro Video}}

\codeexample{language={python}, code={print("Hello")}, output={Hello}}

\practice{
title={Try It},
language={python},
startercode={
print("Hello World")
},
expectedoutput={Hello World}
}

\quiz{
title={Quick Check},
question={What does 2 + 2 equal?},
optionA={3},
optionB={4},
optionC={5},
correct={B},
explanation={Basic arithmetic.}
}

\discussion{prompt={What was the hardest part of this lesson?}}

\resource{type={website}, title={Python Docs}, url={https://docs.python.org}}

\checkpoint{title={Lesson 1 complete!}}

\project{
title={Mini Project},
description={Build something small.},
instructions={Follow the steps in the workspace.},
colab={https://colab.research.google.com}
}

\assignment{
title={Optional Assignment},
instructions={Submit your notes as PDF.}
}

}
}

\end{document}
